\documentclass{report}
\usepackage{amsmath,amssymb}
\setlength{\parindent}{0mm}
\setlength{\parskip}{1em}
\begin{document}
\begin{center}
\rule{6in}{1pt} \
{\large Xiangmin Jiao \\
{\bf A Multigrid Generalized Finite Difference (GFD) Method}}

Dept of Applied Math \& Stat  \\ Stony Brook University \\ Stony Brook \\ NY 11794
\\
{\tt xiangmin.jiao@stonybrook.edu}\\
Cao Lu\end{center}

\renewcommand{\vec}[1]{\mbox{\boldmath $#1$}}

Finite differences are well-established numerical methods for solving
differential equations, but they were limited to structured meshes.
The \emph{generalized finite differences} (\emph{GFDs}) extend the
classical finite differences to unstructured meshes with high-order
accuracy. In this paper, we propose a GFD method for elliptic PDEs,
and tightly integrate it with a multigrid solver. We refer to this
integrated approach as \emph{multigrid GFD}. Specifically, our method
discretizes the PDE and the boundary conditions using GFD to obtain
a sparse linear system
\begin{equation}
\vec{A}\vec{u}=\vec{b},\label{eq:linearsys}
\end{equation}
where $\vec{A}$ is in general asymmetric and may be singular for
PDEs with Neumann or periodic boundary conditions. We solve the resulting
linear system using a hybrid geometric+algebraic method, or \emph{HyGA}.
The core of \emph{HyGA} is a geometric multigrid method, which rediscretizes
the PDE using GFD over a series of hierarchical meshes, combined with
an algebraic multigrid at the coarsest level.

We develop GFD methods based on a \emph{weighted least squares} (\emph{WLS})
formulation locally over a weighted stencil at each point, which generalizes
the classical interpolation-based finite differences. We have shown
recently that the WLS can deliver the same order of accuracy as the
interpolation-based approximations, while delivering superior stability,
flexibility, and robustness for multivariate approximations over irregular
unstructured meshes. The GFD method for PDEs is formulated as follows.
Consider a general form of linear, time-independent partial differential
equations
\begin{equation}
\mathcal{P}\, u(\vec{x})=f(\vec{x})\label{eq:linearPDE}
\end{equation}
with some Dirichlet or Neumman boundary conditions, where $\mathcal{P}$
is a linear differential operator (such as the Laplacian operator)
and $\vec{x}\in\mathbb{R}^{d}$ for $d\ge1$. Let $\Psi(\vec{x})=\{\psi_{j}(\vec{x})\}$
denote the set of Dirac delta functions at the vertices of a mesh.
We obtain a linear equation for each $\psi_{j}$ as
\begin{equation}
\int_{\Omega}\mathcal{P}\, u(\vec{x})\psi_{j}(\vec{x})\,
d\vec{x}=\int_{\Omega}f(\vec{x})\psi_{j}(\vec{x})\,
d\vec{x},\label{eq:weak_form}
\end{equation}
and then the boundary conditions may modify the linear system. We
then introduce a set of basis functions $\Phi(\vec{x})=\{\phi_{i}(\vec{x})\}$
to approximate $\mathcal{P}u$ as $\mathcal{P}\,
u(\vec{x}_{0})=\vec{u}^{T}\mathcal{P}\,\Phi(\vec{x}_{0})\approx\vec{u}^{T}\vec{d}$,
where $\vec{d}$ is obtained from weighted least squares over the
stencil around point $\vec{x}_{0}$. Suppose $\Phi$ and $\Psi$ are
both stored in column vectors. Then in (\ref{eq:linearsys})
\begin{equation}
\vec{A}=\int_{\Omega}\Psi(\vec{x})\left(\mathcal{P}\,\Phi(\vec{x})\right)^{T}\,
d\vec{x}\label{eq:coeff_mat}
\end{equation}
and $\vec{b}=\int_{\Omega}\Psi(\vec{x})f(\vec{x})\, d\vec{x}$.

To construct the geometric multigrid GFD, which is the core of our
multigrid GFD approach, we derive the restriction and prolongation
operators rigorously for hierarchical meshes. Specifically, let $\Phi^{(k)}$
denote the set of basis functions on the $k$th finest mesh, and let
$\Psi^{(k)}$ denote the Dirac delta functions on the $k$th finest
mesh. Let $\vec{A}^{(k)}$ denote the coefficient matrix at the $k$th
level, i.e.,
$\vec{A}^{(k)}=\int_{\Omega}\Psi^{(k)}(\vec{x})\left(\mathcal{P}\,\Phi^{(k)}(\vec{x})\right)^{T}\,
d\vec{x}.$
In a two-level setting, let $\vec{R}_{\Phi}^{(1,2)}$ denote a \emph{restriction
matrix} of the functional space such that $\Phi^{(2)}=\vec{R}_{\Phi}^{(1,2)}\Phi^{(1)}$,
where $\vec{R}_{\Phi}^{(1,2)}\in\mathbb{R}^{n_{2}\times n_{1}}$ with
$n_{k}=\vert\Phi^{(k)}\vert$. Similarly, $\Psi^{(2)}=\vec{R}_{\Psi}^{(1,2)}\Psi^{(1)}$
for $\vec{R}_{\Psi}^{(1,2)}\in\mathbb{R}^{n_{2}\times n_{1}}$. We
show that the restriction matrix $\vec{R}$ and the prolongation matrix
$\vec{P}$ in a two-grid method are
\begin{equation}
\vec{R}=\vec{R}_{\Psi}^{(1,2)}\mbox{ and
}\vec{P}=\left(\vec{R}_{\Phi}^{(1,2)}\right)^{T},\label{eq:restrict_prolong}
\end{equation}
so that $\vec{A}^{(2)}=\vec{R}\,\vec{A}^{(1)}\vec{P}$. In other words,
$\vec{R}$ corresponds to an $n_{2}\times n_{1}$ permutation matrix
(i.e., an injection operator) that selects the entries in the residual
vector corresponding to the coarse-level nodes in the fine-level mesh.
$\vec{P}$ is the transpose of of the restriction operator from the
basis functions $\Phi^{(1)}$ to $\Phi^{(2)}$, and it would correspond
to the interpolation of nodal values if $\Phi^{(k)}$ were composed
of Lagrange basis functions. Applying these operators to adjacent
grids in a multilevel method, we then obtain a geometric multigrid
for GFD, where the coefficient matrices at all levels are obtained
directly from GFD discretizations without performing sparse matrix-matrix
multiplications. By coupling this geometric multigrid with algebraic
multigrid at the coarsest level, our proposed HyGA method then combines
the rigor, high accuracy and runtime-and-memory efficiency of geometric
multigrid with the flexibility of algebraic multigrid, and at the
same time it is relatively easy to implement.

In this work, we also investigate the convergence properties of the
multigrid solver for the GFDs of elliptic PDEs. We show the multigrid
solver may fail to converge for high-order GFD discretizations when
using Jacobi or Gauss-Seidel iterations as smoothers. To overcome
this issue, we propose new weighting schemes for WLS to improve the
diagonal dominance of the linear system from the GFD, and also adopt
a damped Gauss-Seidel iteration to enable more robust convergence
of multigrid methods. We present the derivations and the numerical
experimentations of our method for a number of test problems with
various boundary conditions.


\end{document}
